% Generated draft: requires mathematical and source review.
\documentclass[11pt]{article}
\usepackage[margin=25mm]{geometry}
\usepackage{fontspec}
\setmainfont{Latin Modern Roman}
\usepackage{amsmath,amssymb,mathtools,unicode-math}
\setmathfont{Latin Modern Math}
\usepackage{enumitem,needspace,xurl,hyperref}
\setlength{\parindent}{0pt}
\setlength{\parskip}{5pt}
\setlength{\emergencystretch}{3em}
\newcommand{\N}{\mathbb{N}}
\newcommand{\Z}{\mathbb{Z}}
\newcommand{\Q}{\mathbb{Q}}
\newcommand{\R}{\mathbb{R}}
\newcommand{\C}{\mathbb{C}}
\newcommand{\F}{\mathbb{F}}
\newcommand{\A}{\mathbb{A}}
\newcommand{\PP}{\mathbb{P}}
\newcommand{\E}{\mathbb{E}}
\newcommand{\eps}{\varepsilon}
\newcommand{\dd}{\,\mathrm{d}}
\begin{document}
\section*{UnsolvedMath ID: 20003154}
\textbf{Source problem code:} catalogue entry
\subsection{Definitions and mathematical statement}
A Riemannian submersion $\pi:(G,g)\to(B,h)$ is a smooth surjection whose differential restricts to an isometry from the orthogonal complement of each fiber tangent space onto $TB$. The source asks for a classification when $g$ is bi-invariant and explicitly notes that not every such map is a biquotient submersion. It does not specify compactness, connectedness, fiber conditions, or classification equivalence. A biquotient submersion here is the quotient by a free isometric left--right action $g\mapsto agb^{-1}$ of an appropriate subgroup of $G\times G$.

The supplied partial classification assumes $G$ compact, connected, and semisimple, with connected oriented one-dimensional totally geodesic fibers. The vertical orientation gives a global unit field $V$. The aid asserts that the submersion metric together with total geodesicity makes $V$ Killing, and hence
\[
V=X^L+Y^R,\qquad X^L(g)=\dd L_g(X),\quad Y^R(g)=\dd R_g(Y),
\]
for $X,Y\in\mathfrak g$. Write the orthogonal simple-ideal decomposition $\mathfrak g=\bigoplus_i\mathfrak g_i$, $X=\sum_iX_i$, $Y=\sum_iY_i$. Its factor-support criterion is
\[
\|X^L+Y^R\|\text{ constant on }G
\quad\Longleftrightarrow\quad
\forall i,\ X_i=0\text{ or }Y_i=0.
\]
The left-trivialized cross term is $\langle X,\operatorname{Ad}_{g^{-1}}Y\rangle$; independent variation in each simple factor gives the criterion. Semisimplicity is essential to this particular statement; central torus directions are not covered by the same either-or rule.

The flow has the convention
\[
\varphi_t(g)=e^{tY}ge^{tX}.
\]
In the submersion setting the fibers are compact connected geodesic circles, and the globally oriented unit flow has a common fiber period, producing a free effective isometric circle action after dividing by its global period. Thus the submersion is metrically equivalent to this orbit projection, with the base carrying its quotient metric. On a simple domain one of $X,Y$ vanishes, so the map is equivalent to a one-sided circle-coset projection. Metric equivalence here allows source and target isometries intertwining the submersions.

The support condition alone only classifies constant-length fields; arbitrary such one-parameter flows need not be periodic. Closedness, period lattice, and freeness must still be checked for a global circle quotient. For finite central quotients of products of simply connected simple groups, the Lie-algebra calculation lifts, but periods, effective circle parameters, fiber lifting, and descent must be checked on the actual quotient. They are not determined solely by a list of nonzero simple-ideal components.

The aid reports that the unrestricted problem remains open and includes known examples beyond action quotients, but supplies neither their formulas nor full references. No classification is asserted after dropping total geodesicity, orientation, fiber dimension one, connectedness, compactness, or semisimplicity.
\subsection{Short English statement}
The supplied theorem treats a narrow regime: oriented geodesic-circle fibers on a compact connected semisimple bi-invariant group. Their unit field is reported Killing, and constant length forces its left and right components to occupy disjoint simple factors. The given fibers then supply the periodic free action. A support condition by itself does not supply periodicity or handle finite-central-quotient data, and the full source classification remains broader.
\subsection{Sources}
\begin{itemize}
\item Local registry: permanent catalogue entry, exact source code catalogue entry; \texttt{aim-topology-notes.json}, \texttt{notes[241]}.
\item AIM workshop \emph{Manifolds with nonnegative sectional curvature}, Problem 29, \url{https://aimath.org/WWN/nnsectcurvature/nnsectcurvature.pdf}. The source records the workshop observation that submersions need not be biquotient submersions. No proposer or proposal year is supplied; original remarks and literature field are empty.
\item Verification reports agreement with the source and no OCR repair, and explicitly notes the missing compactness, connectedness, fiber, and equivalence hypotheses.
\item Supplied research attempt 1, registry updated 14 August 2026: partial result, low novelty confidence, machine-generated and not peer reviewed. It supplies the one-dimensional totally geodesic factor-support theorem and finite-central-quotient caveat, but no complete authoritative references for this theorem or the reported general non-action examples.
\end{itemize}
\end{document}
